% In this file you should put all LaTeX macros and settings to be used both by
% the pdf version and the web version.
% This should be most of your macros.

% The theorem-like environments defined below are those that appear by default
% in the dependency graph. See the README of leanblueprint if you need help to
% customize this. 
% The configuration below use the theorem counter for all those environments
% (this is what the [theorem] arguments mean) and never resets it.
% If you want for instance to number them within chapters then you can add
% [chapter] at the end of the next line.
\newtheorem{theorem}{Theorem}
\newtheorem{proposition}[theorem]{Proposition}
\newtheorem{lemma}[theorem]{Lemma}
\newtheorem{corollary}[theorem]{Corollary}
\newtheorem{conjecture}[theorem]{Conjecture}

\theoremstyle{definition}
\newtheorem{definition}[theorem]{Definition}

\newcommand{\C}{\mathbb{C}}
\newcommand{\PP}{\mathbb{P}}
\newcommand{\K}{\mathrm{K}}
\newcommand{\Sym}{\mathrm{Sym}}
\newcommand{\Pic}{\mathrm{Pic}}
\newcommand{\Gr}{\mathrm{Gr}}
\newcommand{\cO}{\mathcal{O}}
\newcommand{\cM}{\mathcal{M}}
\newcommand{\cS}{\mathcal{S}}
\newcommand{\cG}{\mathcal{G}}
\newcommand{\cZ}{\mathcal{Z}}

% MR4433079: Intersection complexes and unramified L-factors — project macros
\newcommand{\cY}{\mathcal{Y}}
\newcommand{\cA}{\mathcal{A}}
\newcommand{\cD}{\mathcal{D}}
\newcommand{\cB}{\mathcal{B}}
\newcommand{\sY}{\mathscr{Y}}
\newcommand{\sM}{\mathscr{M}}
\newcommand{\sA}{\mathscr{A}}
\newcommand{\checkLambda}{\check{\Lambda}}
\newcommand{\checkG}{\check{G}}
\newcommand{\checkT}{\check{T}}
\newcommand{\checkfg}{\check{\mathfrak{g}}}
\newcommand{\checknu}{\check{\nu}}
\newcommand{\checklambda}{\check{\lambda}}
\newcommand{\checktheta}{\check{\theta}}
\newcommand{\mfc}{\mathfrak{c}}
\newcommand{\mfB}{\mathfrak{B}}
\newcommand{\IC}{\mathrm{IC}}
\newcommand{\Bun}{\mathrm{Bun}}
\newcommand{\LocSys}{\mathrm{LocSys}}
\newcommand{\Maps}{\mathrm{Maps}}
\newcommand{\Fr}{\mathrm{Fr}}
\newcommand{\mfo}{\mathfrak{o}}
\newcommand{\GIT}{\!\sslash\!}
\providecommand{\sslash}{/\!\!/}
\newcommand{\ol}{\overline}
\newcommand{\Prim}{\mathrm{Prim}}
\newcommand{\len}{\mathrm{len}}
\newcommand{\msf}[1]{\mathsf{#1}}
\newcommand{\Lplus}{\mathsf{L}^+}
